% Discussion of the simulation results. Macros from results.tex.
\section{Discussion}

\subsection{Static transfer}
At the conditions of the test board the schematic converter is ideal within the resolution of the
bisection: DNL and INL stay within \numrange{\DnlMin}{\DnlMax} LSB and \numrange{\InlMin}{\InlMax}
LSB, while the transition levels are known to \SI{\Resolution}{mV}, about \num{0.016} LSB. There are
no missing codes, the offset is \num{\OffsetLsb} LSB and the end-point LSB is \SI{\Lsb}{mV} instead of
\SI{\LsbIdeal}{mV} (gain error \SI{\GainErr}{\percent}). A least-squares fit of the transition levels to
binary weights gives
\begin{center}\small
\begin{tabular}{lccccc}
\toprule
Bit & 0 & 1 & 2 & 3 & 4 \\
\midrule
$V_{DD} = \SI{\Vdd}{V}$ (board) & \BitWA & \BitWB & \BitWC & \BitWD & \BitWE \\
\ifdefined\RefVdd
$V_{DD} = \SI{\RefVdd}{V}$ (design testbench) & \RefBitWA & \RefBitWB & \RefBitWC & \RefBitWD & \RefBitWE \\
\fi
Ideal & 1 & 2 & 4 & 8 & 16 \\
\bottomrule
\end{tabular}
\end{center}
\ifdefined\RefVdd
The testbench conditions of the design (\SI{\RefVdd}{V}, LSB \SI{\RefLsb}{mV} against
\SI{\RefLsbIdeal}{mV}) give the same picture: gain error \SI{\RefGainErr}{\percent}, $|$DNL$|$ and
$|$INL$|$ below \RefDnlMax{} and \RefInlMax{} LSB. The supply and the reference levels therefore do not
change the static linearity of the circuit, and the lower supply of the board does not starve the
switches at these reference levels.
\fi

The unit capacitor is a MOS gate capacitance that changes by about
\SI{10}{\percent} over the bias range (\SI{\CapMin}{fF} to \SI{\CapMax}{fF}), yet it does not show up in
the static transfer. All units are the same device with the same bias history, and the comparator
always decides with the two top plates close to each other, so the ratios of the DAC capacitances, which
set the weights, do not depend on the input. The small gain error, the same at both operating points, is
a ratio effect as well (the capacitance of the switches and of the dummy units on the top plates).

\subsection{Consequences for the chip measurement}
The measured chip (report of the MSO measurement in \texttt{RISCVConsole/software/mso}) has bit
weights that differ from binary by \SIrange{-8}{7}{\percent}, an INL of about \num{0.5} LSB rms and a
gain error of about \SI{6}{\percent}. None of this is present in the schematic simulation at the same
supply, references and clock.
\IfFileExists{results_pex.tex}{%
The post-layout simulation of the taped-out layout (Section~\ref{sec:postlayout}) reproduces the sign
of all five weight errors, about two thirds of the INL and half of the loss of ENOB: the main cause is
the wiring capacitance of the DAC plates, which differs from bit to bit (the bridges of bits 1 and 3,
the coupling of the MSB plate to its own top plate and of the other floating plates to the opposite top
plate) and between the two halves of the taped-out layout. The rest can come from}{%
Possible causes, outside of what is simulated here:}
\begin{itemize}
  \item the references of the board, which have no decoupling at the chip and are driven by
        \SI{1}{MHz} op-amps: how far they recover within one bit time depends on the bits already decided;
  \item random mismatch of the unit capacitors (not simulated);
  \item the wiring resistance, which the extraction leaves out.
\end{itemize}
A measurement with the references decoupled at the chip would separate the board from the chip.

\subsection{Dynamic performance}
\IfFileExists{table_sweep.tex}{%
With the ideal differential input the transient simulation gives \SineLowEnob{} bits at
\SI{\SineLowFreq}{kHz} and \SineHighEnob{} bits at \SI{\SineHighFreq}{kHz}, practically the
\SineLowSinadSt~dB that the static transfer alone predicts (SINAD$_{st}$). The sampling network adds
almost no error up to the Nyquist frequency: at \SI{\SineHighFreq}{kHz} the SINAD is
\SineHighSinad~dB and H3 rises from \SI{\SineHtreeLow}{dBc} to about \SI{\SineHtreeHigh}{dBc}, still far
below the quantization error. H2 is a numerical zero: the two halves of the converter are identical and
the inputs exactly opposite, so the even orders cancel. The records at \SI{16.7}{kHz} and \SI{94.9}{kHz}
give identical results, as they must for a converter without dynamic errors: with a coherent record and
a number of cycles prime to the number of samples, both sample the same set of phases of the sine.
\IfFileExists{results_fe.tex}{\input{results_fe.tex}%

With the inputs of the board front-end model (Section~\ref{sec:fe}) the ENOB falls to \FeEnobHigh{} bits
at \SI{\FeFreqHigh}{kHz}. Most of this is the smaller amplitude (\FeAmpHigh{} LSB instead of
\FeAmpLow{}). The common mode that moves with the signal produces H2, but at most
\SI{\FeHtwoMax}{dBc} (at \SI{\FeHtwoMaxFreq}{kHz}): the converter rejects the common-mode ripple well.
The chip measurement shows H2 between \SI{-39}{dBc} and \SI{-24}{dBc} from \SI{300}{kHz} to
\SI{650}{kHz}, much more than this, so the measured H2 is not the response of the ADC to the
common-mode ripple of a linear front end; it is distortion of the front end itself (the op-amps near
their slew rate limit).}{}}{%
The sine records are still running; this subsection is completed with them.}

\subsection{Limits of this simulation}
Schematic netlist for the dynamic part; the post-layout netlists (coupling capacitance, no
resistance) for the static transfer only. Typical corner at \SI{25}{\celsius}, no mismatch and no
noise; the board front end enters only through a linear model of its amplitude and phase.
